% Source: https://artofproblemsolving.com/community/c3h2128870_2017_amc_12a_problem_23
% Stable identity: aops-topic-2128870-q01
Let $p,q,r$ be the three distinct roots of $g$, so $g(x)=(x-p)(x-q)(x-r)$. Each of them is a root of $f$, so by the factor theorem $(x-p)(x-q)(x-r)=g(x)$ divides $f(x)$. Since $f$ is monic of degree $4$, we have $f(x)=g(x)(x-t)$ for some number $t$. Expanding,
\[f(x)=x^4+(a-t)x^3+(1-at)x^2+(10-t)x-10t .\]
Comparing with $f(x)=x^4+x^3+bx^2+100x+c$:
\[10-t=100\Rightarrow t=-90,\qquad a-t=1\Rightarrow a=-89,\qquad b=1-at=1-8010=-8009,\qquad c=-10t=900 .\]
Then
\[f(1)=1+1+b+100+c=102-8009+900=-7007,\]
which also equals $g(1)(1-t)=(1-89+1+10)\cdot91=-77\cdot91$. The answer is $\boxed{\textbf{(C)}\ -7007}$.

(The hypothesis is consistent: $g(x)=x^3-89x^2+x+10$ has $g(-1)=-81<0$, $g(0)=10>0$, $g(1)=-77<0$ and $g(100)>0$, so it has three distinct real roots.)
